\section{积分如何表示连续累积}
\subsection{从分割求和到定积分}
对非均匀细杆，不能用某一点的线密度乘全长来求质量。将杆分成许多小段，在第 $i$ 段内取一点 $\xi_i$，以 $\lambda(\xi_i)\Delta x_i$ 近似该段质量。对于连续线密度，当最大段长趋于零时，这些近似之和趋于同一个结果：
\[
 M=\lim_{\max\Delta x_i\to0}\sum_i\lambda(\xi_i)\Delta x_i
 =\int_0^L\lambda(x)\,\dd x.
\]
这种“分割、求和、取极限”的过程定义了定积分。它不要求各段一样长，也不要求函数恒定；连续性保证分割充分细时，段内变化导致的总误差趋零。

一般记 $\int_a^b f(x)\dd x$。$x$ 是积分变量，决定微元如何排列；$\dd x$ 表示所选分割方向的小长度。若自变量是时间，$I(t)\dd t$ 是通过的微小电量；若自变量是体积，$\rho\,\dd V$ 是微小质量。先说明要累加什么，再写出微元，通常比直接寻找可套用的积分公式更可靠。

定积分有方向。约定 $\int_b^a f\,\dd x=-\int_a^b f\,\dd x$，且中间点处可拆分相加。若 $f$ 在横轴两侧变号，积分给出正、负区域的代数和。几何面积应积 $|f|$；路程应积 $|v|$；电流积分则给出净通过电量。能否直接把积分叫作面积，取决于被积函数的符号。

\fig{c2k-sum}{用窄条的高度乘厚度近似局部质量}{\begin{tikzpicture}[x=1cm,y=.9cm]
\draw[axis](0,0)--(4.7,0)node[right]{$x$};
\draw[axis](0,0)--(0,3.2)node[above]{$\lambda$};
\foreach \a/\b/\c in {.8/1.2/.45,1.6/2/.786,2.4/2.8/1.314,3.2/3.6/2.034}
{\draw[fill=blue!12](\a,0)rectangle(\b,\c);}
\draw[world,domain=0:4.1,samples=40]plot(\x,{.3+.15*\x*\x});
\draw[guide](1,0)--(1,.45);
\draw[dim](.8,-.35)--(1.2,-.35);\node[below]at(1,-.35){$\Delta x_i$};
\node[above]at(1,1){$\lambda(\xi_i)$};\draw[guide](1,.9)--(1,.45);
\node[below]at(4,0){$L$};
\end{tikzpicture}}
图 \ref{c2k-sum} 中，每个矩形的横宽对应一小段杆长，高度是段内选定点的线密度。矩形的面积在这里具有质量单位；它只是线密度图下的面积，不是杆在空间中的横截面积。分割细化时，矩形并不是越来越高，而是用更细的横宽跟随密度变化。先明确纵轴的物理意义，才能判断积分得到的是质量、功还是电量。

若累积的是带符号的量，正负部分应保留方向。例如速度图在时间轴下方的部分表示反向位移，其积分为负；若求路程，就要把这一段的符号反过来后再相加。把积分对象从位移改成路程，改变的是微元 $v\,\dd t$ 与 $|v|\,\dd t$，并不是最后对总位移取一次绝对值。


图\ref{c2deep:signed}把这个区别画成一段先向前、后退回的运动。时间轴上方的面积为 $S_+$，下方几何面积的大小为 $S_-$。位移记录位置的净变化，所以反向部分应扣除；路程记录实际走过的长度，所以两部分都要累加：
\[
 \Delta x=S_+-S_-,\qquad s=S_++S_-.
\]
若两块面积相等，物体返回原位，但走过的路程仍为 $2S_+$。这也解释了为什么 $|\int v\dd t|$ 通常不等于 $\int|v|\dd t$：前者在相加时已经把正负位移抵消。
\fig{c2deep:signed}{同一速度图，位移取代数面积，路程取各块面积之和}{\begin{tikzpicture}
\fill[blue!15](0,0)--(0,1.8)--(2,1.8)--(3.2,0)--cycle;
\fill[red!15](3.2,0)--(4,-1.2)--(5.5,-1.2)--(5.5,0)--cycle;
\draw[axis](0,0)--(6,0)node[right]{$t$};\draw[axis](0,-1.7)--(0,2.3)node[above]{$v$};
\draw[thick](0,1.8)--(2,1.8)--(4,-1.2)--(5.5,-1.2);
\node at(1.5,.8){$S_+$};\node at(4.6,-.6){$S_-$};\node[above right]at(3.2,0){转向};
\end{tikzpicture}}


\subsection{微元的近似为什么能够相加}
在位置 $x$ 处，有限小段的准确质量为
$\Delta m=\lambda(x)\Delta x+o(\Delta x)$。
当分割细化，使用 $\dd m=\lambda(x)\dd x$ 保留一阶部分，所有小段的高阶误差总和趋于零。若所取“微元”内部仍有同阶变化，例如把大圆环当成等高小圆柱，却漏掉弧面倾斜因子，误差就不会随分割消失。

这也是画微元图的作用：图中必须同时指明位置、厚度、截面与排列范围。三维物体按平行截面切片时 $\dd V=A(z)\dd z$；求曲面面积时，沿母线的宽度通常是 $\dd s$ 而非 $\dd z$，二者相差倾斜系数。

误差之和趋零还需要一致控制。对长度为 $L$ 的杆，若每一小段内线密度相对所选高度的最大偏差不超过 $\omega$，总质量误差就不超过 $\sum_i\omega\Delta x_i=\omega L$。连续密度在有限闭区间上一致连续，因此分段足够细时，可以使所有小段共用的 $\omega$ 任意小。这一步把“每一段很准”与“全部相加仍然很准”连接起来。
\fig{c2k-surface}{球面窄带的母线宽度与竖直厚度不同}{\begin{tikzpicture}
\draw[axis](0,0)--(2.8,0)node[right]{$r$};\draw[axis](0,0)--(0,2.6)node[above]{$z$};
\draw[world](2,0)arc(0:90:2);
\coordinate(P)at(1.6,1.2);\coordinate(Q)at(1.24,1.68);
\draw[very thick,blue](P)--(Q)node[midway,left]{$\dd s$};
\draw[guide](P)--(1.24,1.2)--(Q);
\draw[dim](1.95,1.2)--(1.95,1.68)node[midway,right]{$|\dd z|$};
\draw[guide](0,0)--(P)node[midway,below]{$R$};
\fill(P)circle(1.5pt)node[below right]{$P$};\node at(4,1){$\dd S=2\pi r\,\dd s$};
\end{tikzpicture}}
图\ref{c2k-surface}的蓝段是球面母线的局部切线段。旋转一周时，它扫出窄带，所以面积应为圆周长乘$\dd s$。在$r>0$的球面部分，固定$R$，对$r(z)^2+z^2=R^2$求导：
\[
2r\frac{\dd r}{\dd z}+2z=0,
\qquad \frac{\dd r}{\dd z}=-\frac zr.
\]
由局部切线段的勾股关系，再把$r^2+z^2=R^2$代回，得到
\begin{align*}
\dd s&=\sqrt{1+\left(\frac{\dd r}{\dd z}\right)^2}\,|\dd z|
=\sqrt{\frac{r^2+z^2}{r^2}}\,|\dd z|
=\frac Rr|\dd z|,\\
\dd S&=2\pi r\,\dd s=2\pi r\frac Rr|\dd z|=2\pi R|\dd z|.
\end{align*}
两极处的表达式可由邻近窄带的极限理解。若直接用$2\pi r|\dd z|$，就漏掉了图中的倾斜因素；即使分割无限细，这个同阶误差也不会消失。

\fig{c2k-band}{窄带周长随极角变化，沿面宽度取弧长}{\begin{tikzpicture}
\draw[world](0,0)circle(1.7);\draw[axis](0,-2)--(0,2.2)node[above]{$z$};
\draw[blue](-1.472,.85)--(1.472,.85);
\draw[guide](0,0)--(1.472,.85);\draw[guide](0,0)--(1.598,.581);
\draw[very thick,red](20:1.7)arc(20:30:1.7);
\node[right]at(1.9,.7){$R\dd\theta$};\draw[guide](1.66,.71)--(1.9,.7);
\draw[dim](-2.1,.581)--(-2.1,.85)node[left]{$|\dd z|$};
\node at(.3,-2.5){图中红弧表示沿面宽度};
\end{tikzpicture}}

图\ref{c2k-band}中沿面宽度为 $R\dd\theta$，水平圆周半径为 $R\sin\theta$；两个因子相乘才得到球面窄带面积。将高度差 $|\dd z|=R\sin\theta\dd\theta$代入后，面积恰为 $2\pi R|\dd z|$，说明等高窄带面积相同，但并不表示各带周长相同。




\section{原函数与微积分基本定理}
\subsection{积分上限函数的变化率}
设 $f$ 在区间内连续，定义 $A(x)=\int_a^x f(s)\dd s$。自变量从 $x$ 到 $x+h$ 时，只多出一条窄带，
\[
 \frac{A(x+h)-A(x)}h=\frac1h\int_x^{x+h}f(s)\dd s\longrightarrow f(x).
\]
所以 $A'(x)=f(x)$。微小区间内的平均高度趋近局部高度，正是“累积”与“求变化率”互相联系的原因。

如果$F'(x)=f(x)$，那么$(A-F)'=f-f=0$，故在一个连通区间内$A(x)-F(x)=C$。代入起点，$C=A(a)-F(a)=-F(a)$，所以$A(x)=F(x)-F(a)$；最后取$x=b$，得到
\begin{equation}\label{eq:c2fundamental}
 \boxed{\int_a^b f(x)\dd x=F(b)-F(a)}.
\end{equation}
不定积分 $\int f(x)\dd x=F(x)+C$ 表示一族原函数，而定积分是给定区间内的一个数。式\eqref{eq:c2fundamental}使积分常数在两端相减时抵消；由初始条件求运动函数时，常数则承载初始位置或初始速度，不能省略。

\fig{c2k-upper}{积分上限移动时只增添边界旁的一条窄带}{\begin{tikzpicture}[x=1cm,y=.9cm]
\fill[blue!8](.4,0)--(2.5,0)--(2.5,1.5)--(.4,1.08)--cycle;
\fill[blue!30](2.5,0)--(3.2,0)--(3.2,1.64)--(2.5,1.5)--cycle;
\draw[axis](0,0)--(4,0)node[right]{$s$};\draw[axis](0,0)--(0,2.4)node[above]{$f(s)$};
\draw[world](.25,1.05)--(3.6,1.72);
\node[below]at(.4,0){$a$};\node[below]at(2.5,0){$x$};\node[below]at(3.2,0){$x+h$};
\draw[guide](2.5,0)--(2.5,1.5);\draw[guide](3.2,0)--(3.2,1.64);
\node at(1.4,.55){$A(x)$};\node[above]at(2.9,1.8){$\Delta A$};
\end{tikzpicture}}
图 \ref{c2k-upper} 中的深色部分是新、旧积分之差。除以横宽 $h$，得到窄带上的平均函数值；当 $h$ 趋于零，连续性使平均值趋于边界值 $f(x)$。这里并没有要求 $f$ 可导，连续就足以得到上限函数可导。若上限本身是 $g(t)$，则边界又以速度 $g'(t)$ 移动，链式法则给出
$\dd[\int_a^{g(t)}f(s)\dd s]/\dd t=f(g(t))g'(t)$。

因此积分常数的处理与求导方向相连。对不定积分，相差常数的函数具有相同导数；对给定初值的运动，初值从这一族函数中选出唯一成员。用定积分表示初值问题，可以在书写上直接保留起始时刻和起始状态，减少漏常数的机会。

\subsection{积分边界移动时的增量}
若累积量为$A(x)=\int_a^x f(s)\dd s$，增加上限$\Delta x$只会增加末端的一小条，故$\Delta A=f(x)\Delta x+o(\Delta x)$，于是$A'=f(x)$。若上下限都变化，$B(t)=\int_{a(t)}^{b(t)}f(x)\dd x$，可先用原函数写成$B(t)=F(b(t))-F(a(t))$，再对两项分别应用链式法则：
\begin{align*}
\dot B
&=F'(b(t))\dot b-F'(a(t))\dot a\\
&=f(b(t))\dot b-f(a(t))\dot a.
\end{align*}
负号由下限移动时被删去的部分决定。例如同一密度分布中一段滑动窗口的质量变化，应比较两端的流入与流出，而不是仅看窗口长度是否不变。若被积函数本身也随时间变化，还须把区间内部的变化加入；这里只讨论固定的 $f(x)$。

这一方法把积分与守恒思想联系起来：总量变化可由边界通量决定。使用之前，应说明积分变量、正向和两端的运动；若把随时间移动的边界误当常量，即使原函数算对，也会漏掉真实的变化率。



图\ref{c2deep:window}将两端同时右移的窗口分成“删去”和“新增”两条窄带。若 $\dot a=\dot b=u>0$，窗口长度不变，内部总量仍以 $u[f(b)-f(a)]$ 的速率变化。右端密度比左端大时，等长的新条带比删去条带含量更多，总量增加；均匀分布才使两项恰好抵消。
\fig{c2deep:window}{固定分布中的移动窗口，右端新增与左端删去分别计入}{\begin{tikzpicture}
\fill[gray!10](1,0)--(4.5,0)--(4.5,2.15)--(1,1.1)--cycle;
\fill[red!25](1,0)--(1.5,0)--(1.5,1.25)--(1,1.1)--cycle;
\fill[blue!25](4.5,0)--(5,0)--(5,2.3)--(4.5,2.15)--cycle;
\draw[axis](0,0)--(6,0)node[right]{$x$};\draw[axis](0,0)--(0,2.9)node[above]{$f$};\draw[thick](.3,.89)--(5.5,2.45);
\foreach \x in {1,1.5,4.5,5}{\draw[guide](\x,0)--(\x,{.8+.3*\x});}
\node[below]at(1,0){$a$};\node[below]at(4.5,0){$b$};\node[red]at(1.1,2){删去};\node[blue]at(4.8,3){新增};
\draw[flow](1,-.7)--(1.5,-.7)node[right]{$u\Delta t$};\draw[flow](4.5,-.7)--(5,-.7);
\end{tikzpicture}}


\subsection{基本积分与适用区间}
求不定积分时，可以先把候选原函数求导，检查是否恰好得到被积式。对幂函数而言，导数会使指数下降一，所以应尝试把指数提高一。由于 $(x^{n+1})'=(n+1)x^n$，还须除以多出的系数，得到
\[
 \int x^n\dd x=\frac{x^{n+1}}{n+1}+C,\qquad n\ne-1.
\]
条件 $n\ne-1$ 不是形式上的附注：当 $n=-1$ 时，候选式的分母为零，而常数 $x^0$ 的导数也不能给出 $x^{-1}$。这时应利用 $(\ln|x|)'=1/x$，所以原函数改为 $\ln|x|+C$。

指数函数的原函数仍是它自身，因为求导后不改变函数形状：$\int\e^x\dd x=\e^x+C$。如果指数中出现系数 $k$，则 $(\e^{kx})'=k\e^{kx}$，积分时必须补上 $1/k$。这个系数补偿与幂函数中的除以 $n+1$ 本质相同，都是把求导产生的额外系数消去。

正弦、余弦的积分还要追踪符号。因为 $(\cos x)'=-\sin x$，$\sin x$ 的一个原函数是 $-\cos x$；因为 $(\sin x)'=\cos x$，$\cos x$ 的一个原函数是 $\sin x$。例如对 $-\cos(kx+\varphi)/k$ 求导时，外面的负号、余弦求导的负号与内层导数 $k$ 共同作用，结果正好为 $\sin(kx+\varphi)$。先这样核验一次，比只记住一个结果更能避免系数和符号错误。

这些原函数关系都要在被积函数连续的区间内使用。$\ln|x|$ 能分别描述零点两侧的原函数，却不能消除被积式在零点的奇异性。因此 $\int_{-1}^{1}\dd x/x$ 不能直接用两端原函数值相减；这类反常积分需分别检查奇点两侧的极限。

在连续区间内，两个原函数的导数相同，它们的差就是常数；如果区间被奇点分成两段，两侧可以具有不同的积分常数。例如 $1/x$ 的正、负半轴分别连通，不能仅凭一个写在式尾的 $C$ 就把零点两侧视为同一个物理过程。

反向求导还能检查参数退化。式 $\int\e^{kx}\dd x=\e^{kx}/k+C$ 只适用于 $k\ne0$；当 $k=0$ 时被积式等于一，原函数为 $x+C$。若需要研究 $k\to0$ 的连续极限，可以选取定积分形式 $(\e^{kx}-1)/k$，其极限才是 $x$。先把初值写清再研究参数极限，能避免被无关的积分常数掩盖实际结果。


图\ref{c2deep:improper}给出不能跨越奇点使用端点差的典型情形。对 $1/x$，若从两侧分别截去零点附近的小区间，则
\[
 \begin{aligned}
 \int_{-1}^{-\varepsilon}\frac{\dd x}{x}
 &=[\ln|x|]_{-1}^{-\varepsilon}
 =\ln\varepsilon-\ln1\longrightarrow-\infty,\\
 \int_{\eta}^{1}\frac{\dd x}{x}
 &=[\ln|x|]_{\eta}^{1}
 =\ln1-\ln\eta\longrightarrow+\infty.
 \end{aligned}
\]
两侧反常积分都不收敛，通常意义下的定积分不存在。人为令 $\eta=\varepsilon$ 得到的对称抵消是另一种限制过程，不能用它代替分别收敛的要求。图上的正、负面积看似对称，不是将两个无穷量相减的依据。
\fig{c2deep:improper}{奇点两侧须分别取极限，不能仅因图形对称就抵消}{\begin{tikzpicture}[x=2cm,y=.55cm]
\draw[axis](-1.5,0)--(1.5,0)node[right]{$x$};\draw[axis](0,-4.8)--(0,4.8)node[above]{$1/x$};
\draw[blue,thick,domain=.23:1.2,samples=70]plot(\x,{1/\x});\draw[red,thick,domain=-1.2:-.23,samples=70]plot(\x,{1/\x});
\draw[guide](.35,0)--(.35,2.857);\draw[guide](-.35,0)--(-.35,-2.857);
\node[below right]at(.35,0){$\eta$};\node[above left]at(-.35,0){$-\varepsilon$};
\draw[flow](.75,-.8)--(.4,-.8);\draw[flow](-.75,.8)--(-.4,.8);
\end{tikzpicture}}


\subsection{换元是同时更换函数、微元与边界}
设 $u=g(x)$，则 $\dd u=g'(x)\dd x$。当被积式含有 $g'(x)\dd x$ 时，
\[
 \int_a^b f(g(x))g'(x)\dd x
 =\int_{g(a)}^{g(b)}f(u)\dd u.
\]
新积分的上下限是新变量的取值，不是旧变量的数值。若采用 $x=\cos\theta$，应同时写出 $\dd x=-\sin\theta\dd\theta$；区间选择还决定 $\sqrt{\sin^2\theta}$ 等于 $\sin\theta$ 还是 $|\sin\theta|$。

有理式可先代数分解。例如
$1/[(x+3)(x+4)]=1/(x+3)-1/(x+4)$，
就化成两个对数积分。任何求出的原函数，都应反向求导核验；求定积分还需检查上下限顺序和结果正负。

换元公式可由链式法则验证。若 $F'=f$，则 $F(g(x))$ 对 $x$ 的导数是 $f(g(x))g'(x)$，对原区间积分后得到 $F(g(b))-F(g(a))$，恰好等于新变量积分的端点差。因此换元中的系数并非凭形式添加，而是负责把同一累积量用新的参数表示。
\fig{c2k-change}{递减换元使新变量的两个端点次序反转}{\begin{tikzpicture}
\draw[axis](-.3,0)--(4.7,0)node[right]{$x$};
\draw[axis](-.3,-1.5)--(4.7,-1.5)node[right]{$u$};
\foreach \x in {0,4}{\draw(\x,-.08)--(\x,.08);\draw(\x,-1.58)--(\x,-1.42);}
\node[above]at(0,0){$a$};\node[above]at(4,0){$b$};
\node[below]at(0,-1.5){$g(b)$};\node[below]at(4,-1.5){$g(a)$};
\draw[flow,blue](0,-.2)--(4,-1.25);
\draw[flow,red](4,-.2)--(0,-1.25);
\end{tikzpicture}}
图 \ref{c2k-change} 表示 $g$ 递减的情形：沿旧变量从 $a$ 走到 $b$，新变量从较大的 $g(a)$ 走到较小的 $g(b)$。保留这一次序时，微元中的负号与方向已经一致；若交换新上下限，还必须再引入一个负号。只有在把区间改写成新变量区间时，才需要特别检查所选分支和单调性。对反三角或根式换元，区间同时决定根号的符号，不能在计算中另行猜测。

\section{面积、体积与质心}
\subsection{连续分布的加权平均}
离散系统质心满足 $x_C=\sum m_ix_i/\sum m_i$。对连续分布取极限，
\[
 x_C=\frac{\int x\,\dd m}{M},\qquad
 y_C=\frac{\int y\,\dd m}{M}.
\]
分子称质量对相应坐标原点的一次矩。密度均匀时可以在分子、分母约去同一个密度；密度变化时不能直接约去。

竖直窄条在 $y=f(x)\ge0$ 与横轴之间，面积微元为 $f(x)\dd x$，该窄条的质心高度为 $f(x)/2$，因此
\[
 y_C=\frac{\tfrac12\int f(x)^2\dd x}{\int f(x)\dd x}.
\]
这里的一次矩可以进一步从窄条内部积分验证。若面密度为 $\sigma$，窄条内高度 $y$ 到 $y+\dd y$ 的质量为 $\sigma\,\dd x\,\dd y$，它对横轴的一次矩为 $y\sigma\,\dd x\,\dd y$。由下到上积分，
\[
 \dd(M y_C)=\sigma\dd x\int_0^{f(x)}y\dd y
 =\sigma\dd x\left[\frac{y^2}{2}\right]_0^{f(x)}
 =\sigma\dd x\left(\frac{f(x)^2}{2}-0\right)
 =\frac{\sigma}{2}f(x)^2\dd x.
\]
再沿 $x$ 累加，并用总质量 $M=\sigma\int f(x)\dd x$ 相除，就得到上式。若改用水平条，必须先写清该高度处左右边界；同一条物理微元不能重复计算两次。

\fig{c2k-centroid}{窄条内部的质心位于上下边界的中间高度}{\begin{tikzpicture}
\draw[axis](0,0)--(4.7,0)node[right]{$x$};\draw[axis](0,0)--(0,2.8)node[above]{$y$};
\draw[world,domain=.25:3.4,samples=30]plot(\x,{2-.15*(\x-1.5)^2});
\draw[fill=blue!18](2.1,0)--(2.3,0)--(2.3,1.904)--(2.1,1.946)--cycle;
\fill(2.2,.962)circle(1.6pt);\draw[flow](4,1.2)--(2.3,.99);
\node[right]at(4,1.2){窄条质心};
\node[above]at(1.3,2.1){$f(x)$};
\draw[dim](1.7,0)--(1.7,1.924)node[midway,left]{$f(x)$};
\draw[guide](1.7,1.924)--(2.2,1.924);
\end{tikzpicture}}
图 \ref{c2k-centroid} 中先用均匀窄条代替一小块平面区域，再把整条质量集中到它自己的质心。该位置只用于计算一次矩，不能把所有窄条都放到曲线顶边。若下边界改为 $g(x)$，窄条面积变成 $[f(x)-g(x)]\dd x$，质心高度为 $[f(x)+g(x)]/2$；相乘后的一次矩为 $\sigma[f(x)^2-g(x)^2]\dd x/2$。这直接说明上下两条边界都怎样进入公式。

对非均匀体，几何对称只有在密度也保留该对称时才能确定质心。例如形状关于纵轴对称但右侧更密，质心仍可能向右偏移。计算结果应位于质量分布的凸包内；若全部质量都在 $0\le x\le L$，却算出 $x_C>L$，应首先检查分子与分母是否使用了同一密度和同一积分范围。

若改用水平条，必须重新由边界解出左、右横坐标。对关于纵轴对称的抛物线薄板 $y=4\ell-x^2/\ell$，高度 $y$处的半宽是 $x_+(y)=\sqrt{\ell(4\ell-y)}$，所以 $\dd A=2x_+(y)\dd y$。同一薄板由水平条与竖直条计算所得的面积、质心应一致；不能改了积分变量却仍保留原条宽。



\subsection{体积分、表面积分与转动惯量}
球体按高度 $z$ 切片，截面是面积 $\pi(R^2-z^2)$ 的圆盘；薄球壳却应按曲面上的窄带求面积。区分实心体与薄壳，是质心题中最常见的建模关口。

转动惯量是另一种加权积分，
\[
 I=\int r_\perp^2\,\dd m.
\]
权重是质量元到转轴距离的平方。对杆绕端点轴，$r_\perp=x$；对圆盘绕中心法线，应取同心圆环，$\dd m=\sigma\,2\pi r\dd r$。距离、密度和微元三者必须来自同一个分割方案。质心一次矩的单位是质量乘长度，转动惯量的单位是质量乘长度平方，二者不可混淆。

微元的选择通常由对称性决定。薄杆沿长度分割，圆盘按同心圆环分割，旋转体按垂直对称轴的圆盘分割，都能使微元内部的密度或轴距便于处理。若一个微元内部的轴距明显不同，就应继续分割，或先在该微元内完成一次积分，不能仅取某个边缘的距离代表全部质量。

还应把物体厚度与选定的数学维数对应起来。实心球用体密度乘体积，薄球壳用面密度乘表面积，细圆环用线密度乘弧长。几何外轮廓相同不代表微元相同。对这些模型分别写出质量元后，总质量、一次矩和转动惯量只是选取不同权重；方法可以迁移，质量元却必须重新核实。

\fig{fv:c2:annulus}{同心圆环的厚度与到转轴的距离具有不同作用}{\begin{tikzpicture}[scale=.9]
\draw[thick](0,0)circle(2);\fill[blue!18,even odd rule](0,0)circle(1.35)(0,0)circle(1.15);
\draw[thick](0,0)circle(1.15);\draw[thick](0,0)circle(1.35);
\draw[->](0,0)--(1.15,0)node[midway,below]{$r$};\draw[->](2.7,1.8)--(1.1,.7)node[pos=0,above]{$\dd r$};
\draw[->](0,0)--(-1.414,-1.414)node[midway,below left]{$R$};
\node at(0,0){$\odot$};\node[below]at(0,-2.3){转轴垂直盘面};
\node[align=left]at(5,0){$\dd m=2\pi\sigma r\dd r$\\[5pt]$\dd I=r^2\dd m$};
\end{tikzpicture}}
图\ref{fv:c2:annulus}中蓝色环带沿半径的宽为 $\dd r$，沿周向却延伸整整一圈。因此面积是一阶量 $2\pi r\dd r$，而不是 $\pi(\dd r)^2$。环带内各点到中心轴的距离在微元近似下均为 $r$，故均匀圆盘有
\[
 I=2\pi\sigma\int_0^Rr^3\dd r
 =2\pi\sigma\left[\frac{r^4}{4}\right]_0^R
 =\frac{\pi\sigma R^4}{2}=\frac12MR^2.
\]
这里的 $r^2$是惯量的权重，$2\pi r$来自微元面积，二者不能相互替代。若转轴改为盘面内的直径，同一圆环各点到新轴的距离不再相同，须先对周向位置再积分，不能继续整环乘 $r^2$。


图\ref{c2k-weights}取同一根非均匀杆上的质量元 $\dd m$，并指定左端转轴 $O$。求总质量时，左右各块只按自身质量相加；求质心一次矩时，同样质量的远端块贡献更大；求转动惯量时，远端块还多乘一次轴距。因而从 $M$ 转到 $Mx_C$、再转到 $I_O$，改变的是所求量的权重，并没有改变杆或密度：
\[
 M=\int\dd m,\qquad Mx_C=\int x\dd m,\qquad I_O=\int x^2\dd m.
\]
\fig{c2k-weights}{同一小段质量，随所求量不同乘以不同位置权重}{\begin{tikzpicture}

\draw[fill=gray!12](0,-.25)rectangle(6,.25);\draw[very thick](0,-.9)--(0,1.2)node[above]{端点轴 $O$};
\draw[fill=blue!30](4.1,-.25)rectangle(4.4,.25);\node[above]at(4.25,.4){$\dd m=\lambda(x)\dd x$};
\draw[dim](0,-.65)--(4.25,-.65)node[midway,below]{$x$};\node[below]at(6,-.25){$L$};
\draw[flow](4.25,-1.15)--(4.25,-1.8);\node at(4.25,-2.2){$\dd m,\quad x\dd m,\quad x^2\dd m$};
\draw[guide](3,-.9)--(3,1);\node[above]at(3,1){中点轴};
\end{tikzpicture}}
若改绕中点轴，图中的物理位置 $x$ 不变，但到新轴的距离是 $|x-L/2|$，惯量应积 $(x-L/2)^2\dd m$。先在图上标出所求量对应的轴或原点，再确定权重，可以避免把“到原点的坐标”与“到转轴的距离”混淆。






\section{积分与运动、做功}
\subsection{沿同一过程累积变化率}
由 $a=\dot v,v=\dot x$ 可依次积分：
\[
 v(t)=v_0+\int_0^t a(s)\dd s,\qquad
 x(t)=x_0+\int_0^t v(s)\dd s.
\]
若加速度给为位置函数，在一个位置单调变化的运动段内，链式法则给出
$a=\dd v/\dd t=(\dd v/\dd x)(\dd x/\dd t)=v\,\dd v/\dd x$。于是定质量的一维牛顿方程可写成$F(x)\dd x=m v\dd v$。位置端点$x_0,x_1$对应速度端点$v_0,v_1$，两边必须沿同一过程积分：
\begin{align*}
\int_{x_0}^{x_1}F(x)\dd x
&=m\int_{v_0}^{v_1}u\dd u
=m\left[\frac{u^2}{2}\right]_{v_0}^{v_1}\\
&=\frac12mv_1^2-\frac12mv_0^2.
\end{align*}
此式是牛顿定律在过程上的积分结果。对转向点，$v=0$ 处应避免无依据地除以 $v$；可以先在单调运动段推导，再以连续性接到转向点。变质量系统的动量流入问题还需另行建立方程，不能直接把这里的 $m$ 换成随时间变化的质量。


\fig{c2deep:work}{同一力随位置变化的图，带方向的窄条面积给出微元功}{\begin{tikzpicture}
\draw[axis](0,0)--(6,0)node[right]{$x$};\draw[axis](0,-1.5)--(0,2.5)node[above]{$F_x$};
\draw[blue,thick](.5,1.8)--(5.5,-1.2);\fill[blue!20](1.8,0)--(2.1,0)--(2.1,.84)--(1.8,1.02)--cycle;
\draw[dim](1.8,-.4)--(2.1,-.4);\node[below]at(1.95,-.4){$\dd x$};\node[above]at(2,1.3){$\dd W=F_x\dd x$};
\node[below]at(.5,0){$x_0$};\node[above]at(5.5,0){$x_1$};\node[below]at(3.5,0){$F_x=0$};
\end{tikzpicture}}

图\ref{c2deep:work}取向右为位移正向。物体右行时，蓝色斜线在横轴上方的面积增加动能，下方的面积减少动能；不能把全部面积大小都算成正功。对一个单调运动段，可由 $m\dot v=F_x$ 两边乘 $v\dd t=\dd x$，得到 $m v\dd v=F_x\dd x$，无需先除以转向点可能为零的速度。若物体在途中转向，应按实际路径及 $\dd x$ 的符号分段累计。


若先由加速度求速度，再由速度求位置，两次积分中的积分变量只是临时记号，应与上限 $t$ 区分。例如 $v(t)=v_0+\int_0^t a(s)\dd s$，这里 $s$ 从起点运行到当前时刻，最终函数仍依赖 $t$。把 $v_0$ 直接写成零，等于加入一个题目未必给出的初始条件；积分方法本身并不允许这样省略。

在力随位置变化的题目中，动能定理先给出 $v^2(x)$，它不能单独决定速度方向。应根据运动支路取 $v=\pm\sqrt{v^2(x)}$，再由 $\dd t=\dd x/v$ 求时间。到达转向点时若平方速度为零，还应检查时间积分是否收敛，以及受力是否会使物体反向。这样，能量式负责联系状态，时间积分负责描述过程，两者承担不同任务。

% thirty-c2-theory
